\documentclass{article}
\usepackage{ctex} % 用于支持中文

\title{数学建模比赛 AI 使用详情说明}
\author{你的名字}
\date{\today}

\begin{document}

\maketitle

\section*{AI 工具版本}
\noindent OpenAI ChatGPT(Nov 5, 2023 version, ChatGPT-4.)

\section*{具体使用目的与环节}
\subsection*{具体目的}
\begin{enumerate}
    \item 思想思路验证
    \item 代码绘图美化
    \item 等等
\end{enumerate}

\subsection*{具体环节}
\begin{enumerate}
    \item 模型建立
    \item 检验
    \item 模型求解等等
\end{enumerate}

\section*{关键交互记录 (重要提示词与回复)}
\subsection*{Query1: <这里写你问的内容>}
\textit{这个部分写你问 GPT 怎么问的，就是你问的 GPT，问题的全部。}
\textit{不管你问多少 GPT 的内容，用 AI 的部分不要涉及到创新，题目思路等核心部分，体现出自己是 AI 使用者不是依赖 AI，第二即使你全文 AI，但是这个部分不要太多，AI 报告总页数不要超过10页内容。3-5 页比较合适。}

\subsection*{Output1: <这里写 AI 的回复>}
\textit{这一部分填上 GPT 怎么回答的。}
\textit{这个是 GPT 回答的全部 (如果太多，就挑选你使用部分+一些其他内容，就是 AI 给你 2000字，你肯定不能全部复制)}

\subsection*{Query2: <这里写你问的内容>}
\textit{这个部分写你问 GPT 怎么问的，就是你问的 GPT，问题的全部。}

\subsection*{Output2: <这里写 AI 的回复>}
\textit{这一部分填上 GPT 怎么回答的。}
\textit{这个是 GPT 回答的全部 (如果太多，就挑选你使用部分+一些其他内容，就是 AI 给你 2000字，你肯定不能全部复制)}


\section*{采纳和人工修改情况}
\subsection*{采纳情况}
\begin{enumerate}
    \item 采纳 1:
    \item 采纳 2:
\end{enumerate}

\subsection*{人工修改}
\begin{enumerate}
    \item 修改 1:
    \item 修改 2:
\end{enumerate}

\end{document}